\documentclass[10pt,a4paper]{article}
\usepackage[utf8]{inputenc}
\usepackage{amsmath}
\usepackage{amsfonts}
\usepackage{amssymb}
\usepackage{enumitem}
\usepackage{amsthm}
\usepackage{verbatim}
\usepackage{bbm}

\newtheorem{example}{Example}
\newtheorem{definition}{Definition}
\newtheorem{theorem}{Theorem}
\begin{document}

\begin{itemize}
	\item discrete time: random walk, markov process
	\item contineous time: brownian motion, stopping times, martingales
	\item  Ito / stochastic calculus
	\item Risk neutral probability measure
	\item Risk neutral valuation, black scholes
	\item Tricks to avoid black scholes, law of unconciousness statisician (=fundamental theorem of asset pricing)
	\item Calll prices 'span' all derivative prices
	\item solving black scholes with neural networks
	\item the problem of trading frequency vs model accuracy
\end{itemize}

\section{Risk neutral evaluation of derivatives}. 
A derivative (portfolio) is a collection of assets. The question is how to price such a derivative. 
Assets 



\bibliographystyle{plain} % You can choose a different style like 'unsrt', 'alpha', etc.
\bibliography{paper}
\end{document}